\documentclass[11pt,letterpaper]{article}
\usepackage[margin=1in]{geometry}
\usepackage{graphicx} % Required for inserting images
\usepackage{amsthm,amsmath,amssymb}
\usepackage{algorithm}
\usepackage{algpseudocode}
\usepackage{xcolor}
\usepackage{soul}
\usepackage{enumitem}

% Theorems
\newtheorem{theorem}{Theorem}[section]
\newtheorem{corollary}[theorem]{Corollary}
\newtheorem{lemma}[theorem]{Lemma}
\newtheorem{claim}[theorem]{Claim}
\theoremstyle{definition}
\newtheorem{definition}[theorem]{Definition}
\newtheorem{remark}[theorem]{Remark}
\newtheorem{observation}[theorem]{Observation}

\newcommand{\reals}{\mathbb{R}}
\newcommand{\Span}{\mathrm{span}}
\newcommand{\coreq}{\mathrm{coreq}}

\newcommand{\calA}{\mathcal{A}}
\newcommand{\calB}{\mathcal{B}}
\newcommand{\calC}{\mathcal{C}}
\newcommand{\calD}{\mathcal{D}}
\newcommand{\calE}{\mathcal{E}}
\newcommand{\calF}{\mathcal{F}}
\newcommand{\calG}{\mathcal{G}}
\newcommand{\calH}{\mathcal{H}}
\newcommand{\calI}{\mathcal{I}}
\newcommand{\calJ}{\mathcal{J}}
\newcommand{\calK}{\mathcal{K}}
\newcommand{\calL}{\mathcal{L}}
\newcommand{\calM}{\mathcal{M}}
\newcommand{\calN}{\mathcal{N}}
\newcommand{\calO}{\mathcal{O}}
\newcommand{\calP}{\mathcal{P}}
\newcommand{\calQ}{\mathcal{Q}}
\newcommand{\calR}{\mathcal{R}}
\newcommand{\calS}{\mathcal{S}}
\newcommand{\calT}{\mathcal{T}}
\newcommand{\calU}{\mathcal{U}}
\newcommand{\calV}{\mathcal{V}}
\newcommand{\calW}{\mathcal{W}}
\newcommand{\calX}{\mathcal{X}}
\newcommand{\calY}{\mathcal{Y}}
\newcommand{\calZ}{\mathcal{Z}}

\newcommand{\cost}{\mathrm{COST}}
\newcommand{\ALG}{\mathrm{ALG}}
\newcommand{\OPT}{\mathrm{OPT}}
\newcommand{\Flow}{\textsc{Flow}}
\newcommand{\ouralg}{\textsc{LineTracker}}
\newcommand{\interval}[1]{\lfloor #1 \rceil}
\newcommand{\off}{\mathrm{off}}
\newcommand{\on}{\mathrm{on}}
\newcommand{\org}{\mathrm{org}}
\newcommand{\rel}{\mathrm{rel}}
\newcommand{\dis}[1]{#1}
\newcommand{\new}{\mathrm{new}}

\newcommand{\solution}{\medskip\noindent\textbf{Solution.}\par\medskip}

\usepackage[hidelinks]{hyperref}
\usepackage{cleveref}

\title{47-835: Network Optimization I \\ Homework 1}
\author{Poon Tze-Yang}
\date{July 2026}


\begin{document}

\maketitle

Students in the OR and ACO programs should answer all questions. For the others, the questions marked * are optional. Those marked ** are optional for all.

\begin{enumerate}
\item A graph is \emph{self-complementary} if it is isomorphic to its complement. Show that every self-complementary graph has $4n$ or $4n+1$ vertices for some integer $n \geq 0$.
\end{enumerate}

\solution
If a graph is self-complementary, its complement must have the same number of edges.
Let $k$ be the number of vertices in this graph.
Then the number of edges in the $k$-clique $K_k$ must be even, in order for the edges to be evenly split between the graph and its complement.
We now show that when $k=4n+2$ or $4n+3$, the number of edges in $K_k$ is odd.
\begin{align*}
    \text{\#edges in }K_{4n+2} = \binom{4n+2}{2} = 8n^2+2n+1, \\
    \text{\#edges in }K_{4n+3} = \binom{4n+3}{2} = 8n^2+10n+3. 
\end{align*}
Both of the above quantities are odd, which completes the proof.

\begin{enumerate}[resume]
\item Consider a bipartite graph $G = (V,E)$ with bipartition $(A,B)$: $V = A \cup B$. Assume that, for some vertex sets $A_1 \subseteq A$ and $B_1 \subseteq B$, there exists a matching $M_A$ covering all vertices in $A_1$ and a matching $M_B$ covering all vertices in $B_1$. Prove that there always exists a matching covering all vertices in $A_1 \cup B_1$.
\end{enumerate}

\solution
Consider $M:=M_A\cup M_B$. 
The degree of each vertex in $M$ is at most $2$: each vertex is incident to at most one edge in $M_A$ and one edge in $M_B$.
Hence, $M$ decomposes into disconnected paths and cycles. 
All vertices in each cycle can be covered by the matching that takes alternating edges of the cycle.
For all paths of odd length, selecting alternating edges starting from an endpoint forms a matching that covers all vertices of that path.
For paths of even length, there are more vertices on the path on one side of the partition than the other.
WLOG let this side be $A$. 
On the $A$ side, at least one of the vertices must not be in $A$ (otherwise there are not enough vertices in the other partition to match to).
If one of the non-endpoint vertices not in $A$, the two incident edges must be from $M_B$, which is a contradiction.
Therefore, one of the endpoints must not be in $A$, which means we can reduce to an odd-length path and cover all required vertices.



\begin{enumerate}[resume]
\item Deduce Hall's theorem from K\"onig's theorem.
\end{enumerate}

\solution
K\"onig's theorem: in any bipartite graph, the maximum number of edges in a matching equals the minimum number of vertices in a vertex cover.

Hall's theorem: a finite bipartite graph has a complete matching of one set of vertices if and only if every subset of that group connects to a neighborhood of at least equal size.

The forward direction of Hall's theorem is straightforward.
Let the two vertices of the bipartite graph be $A\cup B$, with $A$ having the complete matching.
By considering the neighborhood using only edges of the matching, we can immediately see that for any $S\subseteq A$ the size of the neighborhood is at least $|S|$.

It remains to show that if every subset $S\subseteq A$ connects to a neighborhood $N(S)$ such that $|N(S)|\geq |S|$, then there exists a complete matching of $A$.
Using K\"onig's theorem, this is equivalent to showing that the minimum number of vertices in a vertex cover is $|A|$. 
Suppose for contradiction that there was a vertex cover $C$ of size $<|A|$.
Then $A\backslash C \neq \emptyset$, since no proper subset of $A$ is a vertex cover.
Consider all edges incident to $A\backslash C$. 
To cover them, $C\backslash A$ must contain all vertices in $N(A\backslash C)$.
Therefore 
\begin{align*}
    |C| = |C\cap A| + |C\backslash A| \geq |C\cap A| + |N(A\backslash C)|\geq |C\cap A| + |A\backslash C| = |A|,
\end{align*}
where in the second inequality we used K\"onig's theorem. This is the desired contradiction.


\begin{enumerate}[resume]
\item Given a graph $G = (V,E)$, its edge coloring number is the smallest number of colors needed to color the edges in $E$ so that any two edges having a common endpoint have a different color.
\begin{enumerate}
\item Show that the edge coloring number of a bipartite graph $G$ is always equal to its maximum degree $\Delta$ (i.e. the maximum over all vertices $v$ of the number of edges incident to $v$). (Hint: You can prove and use the fact that the edges of any $k$-regular bipartite graph can be partitioned into $k$ matchings).
\item Give an example of a non-bipartite graph for which the edge coloring number is (strictly) greater than $\Delta$.
\end{enumerate}
\end{enumerate}

\solution
\paragraph{(a)} The edge coloring number of every graph is at least the maximum degree, so showing that there always exists a $\Delta$-edge coloring suffices.

First, we prove the hint\footnote{Note, that this is matroid union packing for matching matroid. The packing bound also gives this, with some argument on the min.} by induction on $k$.
Note that for a bipartite graph to be $k$-regular, both sides must have the same number of vertices.
When $k=1$, our graph is a single matching.
When the graph is $(k+1)$-regular, we show that we can find a perfect matching.
This leaves us with a $k$-regular graph, which can be decomposed into $k$ matchings by the inductive hypothesis.
Suppose for contradiction that there is no perfect matching, so by Hall's theorem, there exists a subset of vertices $S$ on one side such that $|N(S)|<|S|$.
But then the total number of edges incident to $N(S)$ is at least $(k+1)|S|$, which contradicts $(k+1)$-regularity by a counting argument, completing the proof of the hint.

Now, given an arbitrary bipartite graph of maximum degree $\Delta$, we show that by adding edges and vertices, we can extend it to a $\Delta$-regular supergraph. 
The hint shows that this supergraph can be decomposed into $\Delta$ matchings, so assigning a different color to each matching suffices to $\Delta$-edge color the original graph.
First, add isolated vertices to equalize the number of vertices on each side of the graph.
Next, for each pair of vertices $u\in A,v\in B$ with degrees smaller than $\Delta$, we add a gadget that increases the degree of each of these vertices by 1.
The gadget comprises a new $K_{\Delta,\Delta}$ graph with one of its edges $(x,y)$ removed, and two edges $(u,y)$ and $(x,v)$ added.
This increases the degrees of $u$ and $v$, and the vertices we added are already $\Delta$-regular.
Repeat this until the graph is $k$-regular, and we can apply the hint as described above. 

(The gadget is used instead of an edge to avoid multi-edges.)

\paragraph{(b)} The triangle graph has an edge coloring number of 3 but a $\Delta$ of 2.

\begin{enumerate}[resume]
\item For an undirected graph $G$, let $\alpha(G)$ denote the maximum cardinality of an independent set of nodes (pairwise non adjacent nodes), and $\alpha'(G)$ denote the maximum cardinality of an independent set of edges (pairwise disjoint edges or matching). Similarly, let $\beta(G)$ denote the minimum cardinality of a vertex cover of $G$ (set of vertices that are incident to all edges), and $\beta'(G)$ denote the minimum cardinality of an edge cover of $G$ (set of edges that are incident to all nodes).
\begin{enumerate}
\item Show that $\alpha(G) + \beta(G) = \alpha'(G) + \beta'(G) = n = |V(G)|$ for all graphs $G$.
\item Show that for a bipartite graph, the size of a maximum independent set is equal to the size of a minimum edge cover.
\end{enumerate}
\end{enumerate}

\solution
\paragraph{(a)}
First, we show that $\alpha(G) + \beta(G) = n$. 
Note that the complement of an independent set is a vertex cover, since for each independent set, each edge in the graph has at least one endpoint that is not in the independent set.
Hence, the complement of the largest independent set must be the smallest vertex cover.

Next, we show that $\alpha'(G) + \beta'(G) = n$. 
Let $C\subset E$ be a minimum cardinality edge cover, and let $M\subseteq C$ be the largest matching in $C$. 
Note that $C$ cannot contain a path of length $\geq 3$, since the middle edge does not cover any additional vertices and can be removed. 
Hence, $C$ comprises disconnected edges and stars, and $M$ takes one edge in each connected component of $C$.
We claim that $M$ is a maximum cardinality matching. 

Suppose for contradiction that there exists a larger matching $M'$.
The symmetric difference $M\Delta M'$ is a graph of degree at most two, so it comprises alternating paths and cycles.
Since $|M'|>|M|$, there must be an alternating path ending in edges of $M'$ at both ends. 
Furthermore, each of the end-vertices were not covered by an edge of $M$, hence they must have been leaves of stars in $C$.
Then, exchanging the edges of $C$ along this alternating path would decrease the size of $C$ while maintaining coverage of all vertices, contradicting the minimality of $C$. 

After removing the vertices covered by $M$, the remaining edges of $C\backslash M$ each cover only one vertex.
Hence, 
\begin{align*}
    n = \text{\#vertices covered by }M + \text{\#remaining edges of }C\backslash M = 2|M| + (|C| - |M|) = |C| + |M|.
\end{align*}

\paragraph{(b)}
By K\"onig's theorem, in bipartite graphs, $\alpha'(G) = \beta(G)$. Hence, by (a), $\alpha(G) = \beta'(G)$.


\begin{enumerate}[resume]
\item We showed that there always exists a solution $x$ to the linear programming formulation of the assignment problem (minimum-cost perfect matching in a bipartite graph) with all components integral.

Reprove this result in the following way. Take a (possibly non-integral) optimum solution $x$. If there are many optimum solutions, take one with as few non-integral values $x_{ij}$ as possible. Show that, if $x$ is not integral, there exists a cycle $C$ with all edges $e = (i,j) \in C$ having a non-integral value $x_{ij}$. Now show how to derive another optimum solution with fewer non-integral values, leading to a contradiction.
\end{enumerate}

\solution

Given a non-integral optimum solution $x$, we first remove all edges $(i,j)$ where $x_{i,j}=0$. 
Note that all vertices incident to edges $(i,j)$ with $x_{i,j}=1$ have degree one, so they are part of a matching, which we can remove from the graph.
We are left with a bipartite graph where every vertex has degree at least two. 
Such a graph must contain at least one cycle, since performing a depth-first walk from any vertex cannot end in a degree-one vertex we have not visited before.

Since the graph is bipartite, the cycle $C$ must be even. 
Partition the odd and even edges of $C=C_1\cup C_2$. 
One of these partitions must have no larger weight than the other, so without loss of generality we have $w(C_1)\leq w(C_2)$.
Let 
\begin{align*}
    \varepsilon := \min\{\min_{e\in C_1}x_e, \min_{e\in C_2}(1-x_e)\},
\end{align*}
such that decreasing all variables in $x$ corresponding to edges in $C_1$ and increasing all variables corresponding to edges in $C_2$ by $\varepsilon$ results a new solution $x'$.
By definition of $\varepsilon$, all variables in $x'$ are in the range $[0,1]$ and at least one more variable is integral comparted to $x$.
Furthermore, the sum of variables for edges incident to each vertex has not changed.
Finally, $w(x')\leq w(x)$, so $x'$ corresponds to an optimal solution with one fewer fractional variable, which is a contradiction.

\begin{enumerate}[resume]
\item[7.*] Consider the following 2-person game on a (not necessarily bipartite) graph $G = (V,E)$. Players 1 and 2 alternate and each selects a (yet unchosen) edge $e$ of the graph so that $e$ together with the previously selected edges form a simple path. The first player unable to select such an edge loses. Show that if $G$ has a perfect matching then player 1 has a winning strategy.
\end{enumerate}

\solution
Player 1 will play a strategy which always chooses edges of the perfect matching. 
This can start from any edge. 
Since matching edges are disconnected, player 2 always plays edges which are not in the matching.

We show that player 1 always has a valid move, which suffices to prove the claim.
Let the edge selected most recent by player 2 be $(u,v)$. 
Let $u$ be the vertex incident to the previous edge, such that $v$ has not previously appeared in the path.
Since $G$ has a perfect matching, $v$ must be incident to an edge $(v,w)$ in the matching.
Furthermore, since players 1 and 2 alternate picking edges in the path, all previous vertices on the path must be incident to some edge picked by player 1.
Therefore, $w$ must be a fresh vertex, and player 1 can pick the edge of the matching incident to $w$.

\begin{enumerate}[resume]
\item[8.*] Consider a bipartite graph $G = (V,E)$ with bipartition $(A,B)$. For $X \subseteq A$, define $\mathrm{def}(X) = |X| - |N(X)|$ where $N(X) = \{b \in B : \exists a \in X \text{ with } (a,b) \in E\}$. Let $\mathrm{defmax} = \max_{X \subseteq A} \mathrm{def}(X)$. Since $\mathrm{def}(\emptyset) = 0$, we have $\mathrm{defmax} \geq 0$.
\begin{enumerate}
\item Generalize Hall's theorem by showing that the maximum size of a matching in a bipartite graph $G$ equals $|A| - \mathrm{defmax}$.
\item For any 2 subsets $X, Y \subseteq A$, show that
\[
\mathrm{def}(X \cup Y) + \mathrm{def}(X \cap Y) \geq \mathrm{def}(X) + \mathrm{def}(Y).
\]
\end{enumerate}
\end{enumerate}

\solution

\paragraph{(a)}
We show both directions for the inequality:

\ul{maximum bipartite matching $\leq |A| - $ defmax:}
Let $X\subseteq A$ be the set which maximizes defmax. 
$X$ can be matched to at most $|N(X)| = |X| - $ defmax vertices. 
$A\backslash X$ can be matched to at most $|A\backslash X|$ vertices. 
Together, the maximum matching contains at most $|A| - $ defmax edges.

\ul{maximum bipartite matching $\geq |A| - $ defmax:}
We prove this statement by induction on defmax. When $\mathrm{defmax}=0$, we simply use Hall's theorem.

Suppose that for graph $G$, $\mathrm{defmax} = k+1$. We show that $G$ has a matching of size $|A| -k-1$.
Let $X\subseteq A$ be a subset with the largest $\mathrm{def}(X)$ and among those, one with minimum size.

\begin{lemma}
    $A\backslash X$ contains a perfect matching with $B\backslash N(X)$.
\end{lemma}
\begin{proof}
    By Hall's theorem, it suffices to show that for each $T\subseteq A\backslash X$, $|N(T)\backslash N(X)|\geq |T|$.
    Suppose for contradiction that there exists a $T\subseteq A\backslash X$ such that $|N(T)\backslash N(X)| < |T|$.
    Then 
    \begin{align*}
        \mathrm{def}(T\cup X) &\geq |T\cup X| - |N(T\cup X)| \\
        &=|T| + |X| - |N(T)\backslash N(X)| - |N(X)| \\
        &> |X| - |N(X)|, 
    \end{align*}
    which contradicts the the definition of defmax of $G$.
\end{proof}

It remains to show that $X$ contains a matching of size $|X|-\mathrm{defmax}$ with $N(X)$. 
Consider any edge $(a,b)$ where $a\in X$ and $b\in N(X)$. 
Since $|X-a|<|X|$, we have that $\mathrm{def}(X-a)< \mathrm{defmax}$. 
Hence, 
\begin{align*}
    & |X-a| - |N(X-a)|<|X| - |N(X)|\\
    &\Rightarrow |N(X-a)|> |N(X)|-1.
\end{align*}
Combined with the fact that $N(X-a)\subseteq N(X)$, we get that $N(X-a) = N(X)$.
But then the graph induced by $(X-a)\cup N(X)$ forms a bipartite graph with $\mathrm{defmax}= k$, so by the inductive hypothesis, there exists a matching of size $|X-a| - k = |X| - \mathrm{defmax}$ from $X-a$ to $N(X)$.
This completes the proof.

\paragraph{(b)}
\begin{align*}
    \mathrm{def}(X\cup Y) + \mathrm{def}(X\cap Y) & = |X\cup Y| - |N(X\cup Y)| + |X\cap Y| - |N(X\cap Y)| \\
    & = (|X\cup Y| + |X\cap Y|) -  (|N(X\cup Y)| + |N(X\cap Y)|)\\
    & = |X| + |Y| -  (|N(X\cup Y)| + |N(X\cap Y)|)\\
    &\geq |X| + |Y| -  (|N(Y)| + |N(X\backslash Y)| + |N(X\cap Y)|)\\
    &\geq |X| + |Y| -  (|N(Y)| + |N(X)|)\\
    & = |X| - |N(X)| + |Y| - |N(Y)| \\
    & = \mathrm{def}(X) + \mathrm{def}(Y),
\end{align*}
where in the first inequality we use that $|N(X\cup Y)|\leq |N(Y)| + |N(X\backslash Y)|$ and in the second inequality we use that $|N(X)|\leq |N(X\backslash Y)| + |N(X\cup Y)|$.

\begin{enumerate}[resume]
\item[9.*]
\begin{enumerate}
\item Is the complement of a disconnected graph always connected? Give a proof or counterexample.
\item Compare the number of unlabeled connected $n$-vertex graphs with the number of unlabeled disconnected $n$-vertex graphs, i.e. which one is larger?
\end{enumerate}
\end{enumerate}

\solution
\paragraph{(a)}
True: The complement of every disconnected graph is connected.

Given an arbitrary disconnected graph $G$ and its complement $\bar{G}$, denote connectivity of $u$ and $v$ in $G$ by $u\leftrightarrow_G v$.
For any two vertices $u,v\in G$, if $(u,v)\notin G$, then $(u,v)\in \bar{G}$ so $u\leftrightarrow_{\bar{G}}v$. 
Otherwise, $(u,v)\in G$, so $u\leftrightarrow_G v$. 
Since $G$ is disconnected, there exists vertex $w$ such that $(u,w),(v,w)\notin G$.
Thus, $(u,w),(v,w)\in \bar{G}$, so $u\leftrightarrow_{\bar{G}}v$.

\paragraph{(b)}
The number of unlabeled connected $n$-vertex graphs is larger. 
From (a), there is an injective mapping from disconnected to connected graphs. 
Furthermore, there exists connected graphs whose complements are also connected graphs (e.g. the $n$-cycle for $n>3$).
Hence the number of unlabeled connected $n$-vertex graphs is larger.

\begin{enumerate}[resume]
\item[10.**] Let $G$ be a bipartite graph with minimum degree $\delta$. Show that the edges of $G$ can be partitioned into $\delta$ edge covers.
\end{enumerate}

\solution
We are given a bipartite graph $G= (L\cup R,E)$.
We construct $\delta$ edge covers and remove their edges from the graph one cover at a time, indexing them $C_0, .., C_{\delta-1}$.
We maintain the invariant that after the edges of the first $i$ covers have been removed, the graph has minimum degree $(\delta-i)$.
This leaves us with some number of leftover edges, which can be arbitrarily added to the existing covers to give a partition of the edges.

Suppose our graph currently has minimum degree $k+1$. 
Each vertex in $L$ will be incident to exactly on edge in the next cover $C$, so the invariant is trivially maintained for these vertices.
For each vertex $v\in R$ with degree $\deg(v)\geq k+1$, it can be incident to at most $\deg(v)-k$ edges in the next cover $C$ for the invariant to be maintainted.

\paragraph{Constructing an edge cover that maintain the invariant.}

We first construct an augmented graph $G'$ by duplicating some vertices of $L$ and $R$. 
Concretely, for each $v\in R$ (resp. $L$), create $\deg(v)-k-1$ copies of it and connect them to vertices of $L$ (resp. $R$) that $v$ is connected to. 
Call the set of vertices we add to the right (resp. left) side $R'$ (resp. $L'$).
\ul{We will show that there is a matching $M$ that covers all vertices of $L\cup R$ in $G'$. }
From such a matching, we can reconstruct an edge cover by merging all copies of each edge in $R$.
By construction, the remaining degrees of each vertex is at least $k$, which preserves the invariant, and this finishes the proof.

In order to construct the desired matching $M$, we first construct a matching $M_1$ from $L$ to $R\cup R'$ that is complete on $L$. 
\begin{lemma}
    A complete matching $M_1$ from $L$ to $R\cup R'$ exists.
\end{lemma}
\begin{proof}
    Suppose such a matching does not exist, then by Hall's theorem there exists some $S\subseteq L$ such that $|S|>|N'(S)|$. (Use $N'$ for $G'$ and $N$ for $G$.)
    Let the number of edges from $S$ to $N(S)$ in $G$ be $\ell$. 
    Then $\ell\geq (k+1)|S|$.
    Let the total number of copies of vertices in $N(S)$ be $c$.
    Then
    \begin{align*}
        c  &= \sum_{v\in N(S)} (\deg(v) - k - 1) \\
        &\geq \ell - (k+1)|N(S)| \\
        &\geq (k+1)(|S|-|N(S)|).
    \end{align*}
\end{proof}


\bibliographystyle{alpha}
\bibliography{refs}

\end{document}
